\documentclass[a4paper,11pt]{jsarticle}
\usepackage[dvipdfmx]{graphicx}
\usepackage{amsmath}
\usepackage{amsfonts}
\usepackage{amssymb}
\usepackage{listings}
\usepackage{xcolor}
\usepackage{geometry}
\usepackage{hyperref}
\usepackage{fancyhdr}
\usepackage{titlesec}
\usepackage{booktabs}
\usepackage{scalebox}
\usepackage{url}

% ページ設定
\geometry{left=25mm,right=25mm,top=30mm,bottom=30mm}

% コードブロックの設定
\lstset{
    basicstyle=\ttfamily\small,
    breaklines=true,
    frame=single,
    numbers=left,
    numberstyle=\tiny,
    keywordstyle=\color{blue},
    commentstyle=\color{green!60!black},
    stringstyle=\color{red},
    backgroundcolor=\color{gray!10},
    showstringspaces=false
}

% タイトル設定
\title{アーセナルベース（ARMS）カードOCRシステムの開発記録}
\author{開発者}
\date{\today}

\begin{document}

\maketitle

\begin{abstract}
この記事では、アーセナルベース（ARMS）というカードゲームのカード情報を自動的に読み取るOCRシステムの開発について記録します。このシステムは、カード画像から構造化されたデータを抽出し、ゲームデータベースの構築を自動化することを目的としています。OpenAI GPT-4oを活用したマルチモーダルAI技術により、従来のOCR技術では困難だった複雑なカード情報の構造化抽出を実現しました。
\end{abstract}

\tableofcontents
\newpage

\section*{はじめに}

この記事では、``アーセナルベース（ARMS）''というカードゲームのカード情報を自動的に読み取るOCRシステムの開発について記録します。このシステムは、カード画像から構造化されたデータを抽出し、ゲームデータベースの構築を自動化することを目的としています。

\section*{プロジェクト概要}

\subsection*{開発背景}
アーセナルベース（ARMS）は、ガンダムシリーズを題材としたカードゲームです。カードにはMS（モビルスーツ）カードとPL（パイロット）カードがあり、それぞれに詳細な能力値やスキル情報が記載されています。これらの情報を手動で入力するのは非常に時間がかかるため、OCR技術を用いた自動化システムの開発に取り組みました。

\subsection*{システムの特徴}
\begin{itemize}
    \item \textbf{多段階OCR処理}: 基本情報の抽出から詳細なスキル分析まで段階的に処理
    \item \textbf{構造化データ出力}: JSON形式で統一されたデータ構造を提供
    \item \textbf{エラーハンドリング}: 処理失敗時の詳細なログ記録
    \item \textbf{拡張性}: 新しいカードタイプやスキルに対応可能な設計
\end{itemize}

脚注は \verb=\footnote{...}= というように表記して、文章を補足することができます\footnote{\url{https://www.gundam-ab.com/}}。

\section*{アーセナルベース（ARMS）のカード種類}

アーセナルベース（ARMS）には、大きく分けて以下のカード種類が存在します：

\subsection*{MSカード（モビルスーツカード）}

MSカードは、ゲームの中心となるモビルスーツを表すカードです。各MSには以下の情報が記載されています：

\subsubsection*{基本情報}
\begin{itemize}
    \item \textbf{カードID}: 一意の識別番号（例：FQ01-001）
    \item \textbf{カード名}: MSの名称（例：ガンダム試作3号機）
    \item \textbf{モデル}: 機体の正式名称（例：RX-78GP03S GUNDAM GP03S）
    \item \textbf{コスト}: デッキ構築時のコスト（1-10）
    \item \textbf{カテゴリ}: 機体の分類（近距離、遠距離、機動）
\end{itemize}

\subsubsection*{能力値}
\begin{itemize}
    \item \textbf{機動力}: 移動力と回避力に影響
    \item \textbf{遠距離攻撃力}: 射撃攻撃の威力
    \item \textbf{近距離攻撃力}: 格闘攻撃の威力
    \item \textbf{HP}: 耐久力
\end{itemize}

\subsubsection*{地形適性}
\begin{itemize}
    \item \textbf{地上}: A/S/B/C/D（Aが最高）
    \item \textbf{宇宙}: A/S/B/C/D
    \item \textbf{砂漠}: A/S/B/C/D
    \item \textbf{水中}: A/S/B/C/D
\end{itemize}

\subsubsection*{武器情報}
\begin{itemize}
    \item \textbf{メイン武器}: 主要攻撃手段（射程、攻撃種別）
    \item \textbf{サブ武器}: 補助攻撃手段（射程、攻撃種別）
\end{itemize}

\subsection*{PLカード（パイロットカード）}

PLカードは、MSを操縦するパイロットを表すカードです。各パイロットには以下の情報が記載されています：

\subsubsection*{基本情報}
\begin{itemize}
    \item \textbf{カードID}: 一意の識別番号（例：PL01-001）
    \item \textbf{パイロット名}: パイロットの名称
    \item \textbf{所属}: 所属組織（地球連邦軍、ジオン公国軍など）
    \item \textbf{コスト}: デッキ構築時のコスト
\end{itemize}

\subsubsection*{スキル情報}
\begin{itemize}
    \item \textbf{USP}: ユニークスキル（パイロット固有の能力）
    \item \textbf{SQ}: スキル（複数のパイロットで共有される能力）
    \item \textbf{リンクアビリティ}: 特定条件での追加効果
\end{itemize}

\section*{技術的実装}

\subsection*{OCRシステムの構成}

このOCRシステムは、以下の主要コンポーネントで構成されています：

\begin{enumerate}
    \item \textbf{画像前処理}: カード画像の正規化と品質向上
    \item \textbf{基本OCR}: カードの基本情報（ID、名前、コストなど）の抽出
    \item \textbf{詳細OCR}: 能力値、武器、スキルなどの詳細情報の抽出
    \item \textbf{SP OCR}: MSカードのSP部分を詳細処理
    \item \textbf{SQ分析}: PLカードのSQ関連スキルを分析
\end{enumerate}

\subsection*{処理フロー}

システムの処理フローを表\ref{processing-flow}に示します。

\begin{table}[htb]
  \caption{OCRシステム処理フロー}
  \label{processing-flow}
  \begin{center}
    \scalebox{0.8}{
      \begin{tabular}{cll}
        \toprule
        段階 & 処理名 & 関数名 \\
        \midrule
        入力 & カード画像URL/ファイル & - \\
        \midrule
        1 & 基本OCR処理 & \texttt{extract\_structured\_data\_from\_url} \\
        2 & MSカード抽出 & \texttt{extract\_ms\_cards\_from\_basic\_results} \\
        3 & SP部分OCR処理 & \texttt{extract\_sp\_details\_from\_image} \\
        4 & PLカードSQ分析 & \texttt{analyze\_pl\_sq\_skills} \\
        \midrule
        出力 & 構造化JSONデータ & - \\
        \bottomrule
      \end{tabular}
    }
    \end{center}
\end{table}

この4段階の処理フローにより、効率的かつ正確なカード情報の抽出を実現しています。各段階でエラーが発生した場合でも、前の段階の結果を保持し、部分的な成功を記録することができます。

\subsection*{LaTeXコマンドの使用例}

テンプレートに従って、LaTeXコマンドを表示する際は \verb=\section*{...}= のように使用します。また、等幅フォントで表示したい場合は \verb|\texttt{monospace}| としてもよいです（\texttt{monospace} のように表示されます）。

URLは \verb=\url{http://www.example.com}= のように入力すると \url{http://www.example.com} のように表示されます。

この段階的処理により、以下のメリットを実現しました：
\begin{itemize}
    \item \textbf{処理時間の短縮}: 必要なカードのみ詳細処理
    \item \textbf{エラーの局所化}: 段階ごとのエラー管理
    \item \textbf{リソース効率}: API使用量の最適化
\end{itemize}

\section*{結果と成果}

\subsection*{処理実績}
\begin{itemize}
    \item \textbf{対象カード数}: 数百枚のカードを処理
    \item \textbf{成功率}: 基本OCRで90\%以上、SP OCRで80\%以上
    \item \textbf{データ品質}: 手動入力と比較して95\%以上の精度
\end{itemize}

\subsection*{出力データ例}

\begin{lstlisting}[language=JSON]
{
  "card_number": "FQ01-001",
  "card_name": "ガンダム試作3号機",
  "type": "MS",
  "cost": 6,
  "stats": {
    "mobility": 320,
    "ranged_attack": 570,
    "melee_attack": 130,
    "hp": 580
  },
  "terrain_compatibility": {
    "ground": "A",
    "space": "S",
    "desert": "C",
    "water": "C"
  },
  "weapon": {
    "main": {
      "name": "フォールディング・バズーカ",
      "range": 4,
      "type": "遠距離"
    },
    "sub": {
      "name": "ビーム・サーベル",
      "range": 1,
      "type": "近距離"
    }
  },
  "ms_ability": {
    "name": "陣略【妨害／機動】",
    "type": "範囲（敵）",
    "range": 4,
    "cost": 4,
    "description": "一定時間、範囲内に入った敵ユニットの「機動力」をダウンさせるエリアを設置する。※上限3箇所まで。"
  },
  "usp": {
    "name": "ニュータイプ覚醒",
    "condition": "HPが30\%以下で、かつ敵ユニットが2体以上存在する場合",
    "effect": "自身の「機動力」「遠距離攻撃力」「近距離攻撃力」を50\%アップ",
    "limitation": "ゲーム中1回のみ発動可能",
    "activation_cost": 0
  },
  "special_attack": {
    "name": "フォールディング・バズーカ連射",
    "type": "normal",
    "target": "単体（敵）",
    "range": 2,
    "power": 3500
  },
  "link_ability": [
    {
      "name": "ガンダム0083",
      "condition": "デッキに3枚以上",
      "effect": "【機動力】小アップ"
    },
    {
      "name": "的確な一撃",
      "condition": "デッキに4枚以上",
      "effect": "【機動力】【遠距離攻撃力】小アップ"
    }
  ]
}
\end{lstlisting}

\section*{今後の展望}

\subsection*{改善点}
\begin{enumerate}
    \item \textbf{精度向上}: より複雑なカードレイアウトへの対応
    \item \textbf{処理速度}: バッチ処理の最適化
    \item \textbf{新カード対応}: 新しいカードタイプへの拡張
    \item \textbf{UI改善}: 処理結果の可視化ツール
\end{enumerate}

\subsection*{応用可能性}
\begin{itemize}
    \item \textbf{他のカードゲーム}: 同様のOCRシステムの他ゲームへの適用
    \item \textbf{リアルタイム処理}: カメラ撮影からの即座の情報抽出
    \item \textbf{データベース連携}: 既存のゲームデータベースとの統合
\end{itemize}

\section*{まとめ}

このOCRシステムの開発を通じて、AI技術を活用した実用的なソリューションの構築について多くの学びを得ました。特に、プロンプトエンジニアリングの重要性、段階的処理の効果、そしてエラーハンドリングの必要性を実感しました。

\begin{quotation}
    ``AI技術を活用したOCRシステムの開発は、従来の画像認識技術の限界を超える可能性を秘めている。特に、カードゲームのような複雑な情報構造を持つコンテンツにおいて、その効果は顕著である。''
\end{quotation}

カードゲームの情報抽出という具体的な課題に対して、最新のAI技術を適用することで、大幅な作業効率化を実現できました。この経験は、今後他のプロジェクトでも活用できる貴重な知見となっています。

ChatGPT APIの活用により、従来のOCR技術では困難だった複雑なカード情報の構造化抽出が可能になりました。特に、カードゲーム特有の情報構造を理解し、適切なJSON形式で出力できる点が、このシステムの大きな特徴となっています。

\vspace{2em}

\begin{center}
\textit{この記事は、アーセナルベース（ARMS）カードOCRシステムの開発記録として執筆されました。技術的な詳細や実装の工夫について、同様のプロジェクトに取り組む方々の参考になれば幸いです。}
\end{center}

\end{document} 